\section{Classification of graph state sums}\label{sec:graph-classification}

The three separation witnesses also determine which graph families
admit polynomial-size computations of their state sums. Recall that
\(Z_{G,q}\) assigns a separate variable to every choice of the
\(q\)-valued labels on the
edges incident to each vertex, as in Equation~\eqref{eq:graph}.
We measure size in terms of the number \(d\) of vertices and the
alphabet size \(q\). A single vertex of degree \(t\) thus supplies
\(q^t\) distinct input variables. In a multigraph every parallel edge
carries its own label and counts separately towards degrees and cuts.

Write \(m(G)\) for the largest order of a connected component and
\(\Delta(G)\) for the maximum degree. The \emph{cutwidth}
\(\operatorname{cw}(G)\) is the minimum, over vertex orders, of the
largest number of edges crossing from a prefix to its complement.
A \emph{tree decomposition} consists of a tree of vertex sets, called
bags. The bags containing a given vertex form a nonempty connected
subtree, and the two ends of each edge belong to a common bag. The
\emph{treewidth} \(\operatorname{tw}(G)\) is the minimum largest bag
size minus one. Requiring the decomposition tree to be a path gives
\emph{pathwidth}, denoted by \(\operatorname{pw}(G)\).

\begin{theorem}[Polynomial-size graph state sums]\label{thm:graph-classification}
Fix a field and a class \(\mathcal G\) of nonempty finite loopless
multigraphs. In each model below, a single polynomial in \(d+q\)
bounds the size of \(Z_{G,q}\) for every \(G\in\mathcal G\) and
every integer \(q\ge2\) if and only if the indicated parameters
are bounded on \(\mathcal G\).
\begin{center}
\begin{tabular}{@{}ll@{}}
\toprule
\textup{Arithmetic model} & \textup{Bounded graph parameters}\\
\midrule
Formula & \(m(G)\) and \(\Delta(G)\)\\
Algebraic branching program & \(\operatorname{cw}(G)\)\\
Circuit & \(\operatorname{tw}(G)\) and \(\Delta(G)\)\\
\bottomrule
\end{tabular}
\end{center}
For a sequence \(G_n\) with \(|V(G_n)|\le n^{O(1)}\), the same criteria,
uniformly in \(n\), characterize polynomial size for \(Z_{G_n,n}\).
In both formulations the upper bounds hold for exact computation, and
the necessary conditions already hold for degree-component border size.
\end{theorem}

\subsection{The degree obstruction}

Two vertices joined by \(t\) parallel edges give a dot product of two
tables of length \(q^t\). This graph has component order two and
treewidth one for every \(t\). The coefficient matrix of the dot
product is the \(q^t\)-by-\(q^t\) identity. The same obstruction occurs
at any vertex \(v\) of degree \(t\). Flatten the coefficient tensor
across the slot at \(v\) and all remaining slots. Each of its \(q^t\)
rows has nonempty support, since every choice of incident labels
extends to a full edge labeling. The supports are disjoint, since the
table indices at the neighbors determine all these labels. The rank is
therefore \(q^t\) over every field.

This rank bound still detects large tables when affine input forms
are free. For a size-\(s\) computation, let \(U\) be the span of the
linear parts of its affine inputs in the variables at \(v\). There are
\(O(s+1)\) such inputs in each of the three models. In the part of the
output linear in those variables, each product of input forms takes
its \(v\)-linear term from exactly one factor. Extracting a coefficient
in the other slots therefore gives a vector in \(U\). The flattening
has rank at most \(\dim U=O(s+1)\). Since this condition is defined
by vanishing minors of the slot tensor, it holds also on the
degree-component closure.

Consequently a polynomial size bound requires bounded degree. For a
graph class take \(q\ge\max\{2,d\}\) arbitrarily large for each \(G\).
For a sequence take \(q=n\). In both cases the bound
\(q^{\Delta(G)}=O(s+1)\) forces a uniform bound on \(\Delta(G)\).
In the branching-program row, cutwidth already controls degree. In
any order the edges incident to \(v\) cross either the cut immediately
before \(v\) or the cut immediately after it, so
\(\Delta(G)\le2\operatorname{cw}(G)\).

\subsection{Exact computations from graph decompositions}

\paragraph{Formulas.}
The state sum factors over connected components. In a component on
at most \(a\) vertices, enumerate its at most \(q^{\Delta(G)a/2}\)
edge labelings. Each labeling contributes one product of local table
entries. Summing these products within each component and multiplying
the results gives a formula of size \(O((d+1)q^{\Delta(G)a/2})\).
An isolated vertex contributes its single variable.

\paragraph{Branching programs.}
Take an order of cutwidth \(k\), and place one layer after each
prefix. A state records the labels of the at most \(k\) edges that
cross the current cut. At the next vertex, the state gives the labels
on its backward edges. Choose labels for its forward edges, read the
resulting table entry as the transition weight, and retain just the
labels that cross the next cut. Every full edge labeling determines
one source--sink path of the required weight, and every such path
comes from a labeling. There are at most \(dq^{k+\Delta(G)}\)
transitions.

\paragraph{Circuits.}
A tree decomposition of width \(k\) yields a circuit of size
\(O((d+1)q^{\Delta(G)(k+1)})\). Contract adjacent bags whenever one
contains the other, keeping the larger bag, and root the remaining
tree. Each nonroot bag contains a vertex absent from its parent.
By connectedness this is the highest occurrence of that vertex,
so the decomposition has at most \(d+1\) bags.

Assign each vertex table to a bag containing the vertex. Replace
the vertex set of each bag by the list of all edges incident to that
set. The new list has at most \(\Delta(G)(k+1)\) entries and includes
every edge label needed by a table assigned there. For each edge,
the bags containing it still form a connected subtree, since the
subtrees of its endpoints intersect. This construction is related to
the characterization of edge contraction complexity by line-graph
treewidth due to Markov and Shi~\cite{MarkovShi}.

Compute a table at each bag, starting at the leaves. For each assignment
on its edge list, multiply the local vertex entries by the child
tables evaluated on the edges common to the child and the bag. Sum
these values over the labels absent from the parent. The resulting
table is indexed by the labels shared with the parent. Connectedness
ensures that every label used on both sides of this boundary occurs
in the intersection, and that each other label is summed out exactly
once. At the root the output is \(Z_{G,q}\). The total number of
child factors and assigned vertex tables over all bags is \(O(d+1)\),
giving the stated gate count.

\subsection{Transferring lower bounds through graph minors}

To use a graph minor \(H\), we must convert a computation for its
host \(G\) into one for \(Z_{H,q}\). Each vertex of \(H\) is represented
by a connected set of host vertices. The construction routes the labels
entering this set to a single vertex, where its table entry is read.
The required host alphabet depends only on the degree of \(H\).

\begin{lemma}[Projection to a minor]\label{lem:graph-minor}
Let \(H\) be a loopless minor of a finite loopless multigraph \(G\),
and set \(b=\max\{1,\lfloor\Delta(H)/2\rfloor\}\).
If \(q^b\le n\), then \(Z_{H,q}\) is a projection of \(Z_{G,n}\)
using only variables, zero and one.
\end{lemma}

\begin{proof}
Choose disjoint connected branch sets for the vertices of \(H\).
For each edge of \(H\), choose a distinct host edge joining its two
branch sets. Use the first \(q\) states on each chosen edge to carry
its label. Within the branch set for a vertex \(v\), choose a spanning
tree and weight each tree vertex by the number of chosen edges meeting
it. The total weight is \(\deg_H(v)\). Root the tree at a weighted
centroid, so every component below the root has at most half the total
weight. To find such a root, move into any neighboring component of
weight greater than half. Each move leaves less than half the weight
behind, so this walk on a finite tree terminates.

Fix an ordering of the chosen edges. On each tree edge from a child
to its parent, encode in one host state the tuple of labels entering
the child's subtree, in the induced order. The tuple has at most
\(\lfloor\deg_H(v)/2\rfloor\) entries, so \(q^b\le n\) makes all its
values available. At a nonroot vertex, set a local entry to one
precisely when the tuple sent to its parent agrees with the tuples
from its children and the labels of the chosen edges meeting the
vertex. At the root these same data specify the entire tuple at \(v\).
Use the corresponding variable of \(Z_{H,q}\) there. This also covers
an isolated vertex of \(H\), whose tuple is empty.

Pin all unused edges to a single state and give unused vertices
weight one on legal inputs. Set inconsistent table entries and unused
encodings to zero. Each labeling of \(H\) extends uniquely to a
surviving host labeling, because every tuple on a tree edge is forced
by the chosen external labels. The corresponding monomial is exactly
the product of the variables read at the roots.
\end{proof}

\begin{proof}[Necessity in Theorem~\ref{thm:graph-classification}]
First consider a sequence \(G_n\) of polynomially bounded order.
The degree obstruction allows us to assume \(\Delta(G_n)\) bounded.
Each of the remaining unbounded parameters forces its own witnesses.

\begin{itemize}
\item If component order is unbounded, arbitrarily long paths occur
as subgraphs. In a graph of bounded degree, balls of any fixed radius
have bounded order.
\item If cutwidth is unbounded, so is pathwidth. Indeed
\[
 \operatorname{pw}(G)\le\operatorname{cw}(G)
 \le\Delta(G)(\operatorname{pw}(G)+1).
\]
For the first inequality, form each path bag from the next vertex in
an order and the processed vertices with an edge to an unprocessed
vertex. For the second, order the vertices by their first occurrence
in a path decomposition. At each cut, every processed vertex with a
future neighbor belongs to the bag in which the next vertex first
occurs. The crossing edges are therefore incident to at most
\(\operatorname{pw}(G)+1\) processed vertices. The excluded-forest
theorem gives complete binary trees of arbitrarily many levels as
minors~\cite{BRST}.
\item If treewidth is unbounded, square grids of arbitrarily large side
length occur as minors by the excluded-grid theorem
\cite{RSgrid,ChekuriChuzhoy}.
\end{itemize}

Treewidth and pathwidth depend only on the underlying simple graph.
The tree and grid minors just obtained are thus also minors of the
original multigraphs. For every fixed witness size, the corresponding
minor occurs along a subsequence with \(n\to\infty\).
Lemma~\ref{lem:graph-minor} retains \(q=n\) for paths and binary trees
and permits \(q=\lfloor\sqrt n\rfloor\) for grids.

Suppose the host border size were polynomial in \(n\).
Lemma~\ref{lem:projection}, applied with host degree \(|V(G_n)|\),
would give an upper bound \(n^C\) for all projected witnesses, with
\(C\) independent of the witness size. But the short-witness lower
bounds grow as \(q^{\Omega(\log d')}\) for formulas on a path of
\(d'\) vertices, \(q^{\Omega(\sqrt h)}\) for ABPs on a binary tree of
\(h\) levels, and \(q^{\Omega(\sqrt r)}\) for circuits on an
\(r\)-by-\(r\) grid (Theorems~\ref{thm:full-flat-graph},
\ref{thm:b} and~\ref{thm:grid}). Choose the witness size large enough
that its exponent in \(n\) exceeds \(C\), even accounting for the
square root in the grid alphabet. Then take \(n\) sufficiently large
along the corresponding subsequence so that the hypotheses of the
short-witness bound hold. This proves necessity for sequences.

Finally, if a graph class has an unbounded parameter in its row, choose
graphs on which that parameter tends to infinity. Give each graph a
distinct index \(n\) at least its order and fill the other indices from
one fixed graph in the class. A uniform polynomial bound in \(d+q\)
would give a polynomial bound in \(n\) for this sequence at \(q=n\),
contrary to the preceding argument.
\end{proof}
